\documentclass[11pt]{article}
\usepackage[letterpaper,margin=0.6in]{geometry}
\usepackage[T1]{fontenc}
\usepackage{lmodern}
\usepackage{microtype}
\usepackage{enumitem}
\usepackage{titlesec}
\usepackage[hidelinks]{hyperref}
\input{glyphtounicode}
\pdfgentounicode=1

\pagestyle{empty}
\setlength{\parindent}{0pt}
\setlength{\parskip}{0pt}
\setlist[itemize]{leftmargin=15pt,label=\textbullet,itemsep=3pt,topsep=3pt,parsep=0pt,partopsep=0pt}
\titleformat{\section}{\large\scshape}{}{0pt}{}[\titlerule]
\titlespacing*{\section}{0pt}{10pt}{5pt}
\urlstyle{same}
\hypersetup{pdftitle={Daniel David -- Extended Resume},pdfauthor={Daniel David}}

\newcommand{\entry}[2]{\textbf{#1}\hfill #2\par}
\newcommand{\project}[2]{\textbf{#1}\par{\small\textit{#2}}\par}

\begin{document}
\begin{center}
    {\Huge\scshape Daniel David}\\[3pt]
    \small +7 (923) 408 72 93 \enspace\textbar\enspace
    \href{mailto:dayveed.d@gmail.com}{\underline{dayveed.d@gmail.com}} \enspace\textbar\enspace
    \href{https://www.linkedin.com/in/dayveedaniel}{\underline{LinkedIn}} \enspace\textbar\enspace
    \href{https://github.com/dayveedaniel}{\underline{github.com/dayveedaniel}}\\[3pt]
    \textbf{Spoken Languages:} English, Russian
\end{center}
\vspace{-5pt}

\section{Education}
\entry{Innopolis University}{Innopolis, Russia}
\textit{MSc, Artificial Intelligence \& Data Engineering}\hfill Aug. 2024 -- Jul. 2026 \textbar\ CGPA: 4.24\par
\vspace{6pt}
\entry{Tomsk Polytechnic University}{Tomsk, Russia}
\textit{BEng, Software Engineering}\hfill Sep. 2019 -- Jul. 2024 \textbar\ CGPA: 3.85\par

\section{Experience}
\entry{PRAAMS --- Senior Software Developer}{Sep. 2022 -- Present}
\begin{itemize}
    \item Led architecture and development of PRAAMS, an AI-powered co-investor app combining conversational AI, stock analysis, and interactive financial visualizations.
    \item Conducted technical interviews and onboarded new team members.
\end{itemize}
\vspace{6pt}
\entry{Bellissimo Pizza --- ML / Data Engineering Consultant}{Jun. 2026 -- Nov. 2026}
\begin{itemize}
    \item Built an end-to-end ML data pipeline for order-demand and price prediction using Apache Airflow, PostgreSQL, and MinIO; automated ingestion, transformation, and training-data preparation.
    \item Contributed to hourly demand forecasting for 50+ restaurants over 24/48/72-hour horizons, incorporating POS and weather features into model training and evaluation.
    \item Automated model retraining on new data and maintained data pipelines to support recurring forecasting workflows.
\end{itemize}
\vspace{6pt}
\entry{Innopolis University --- Student Researcher}{Aug. 2025 -- Sep. 2026}
\begin{itemize}
    \item Implemented and trained ML/RL models in PyTorch; built web-scraping and big-data analysis pipelines using MongoDB, PySpark, and Apache Hadoop.
\end{itemize}

\section{Selected Projects}
\project{nanoTitan --- Open-Source Mixture-of-Experts Training Stack}{Python, PyTorch, SDPA}
\begin{itemize}
    \item Added configurable SDPA / Flash Attention support to the existing multi-head attention module, retaining an eager attention path as a fallback.
    \item Improved the model architecture by integrating fused QKV projections and RMSNorm.
\end{itemize}
\vspace{6pt}
\project{Multilingual Math Reasoning for Nigerian Languages}{Docker, Postgres, Hugging Face, vLLM, MLX, GSM8K}
\begin{itemize}
    \item Built an ML data pipeline translating GSM8K questions and solutions into Igbo, Hausa, and Yoruba for Tiny Aya supervised fine-tuning on math solving and explanation; produced the \href{https://huggingface.co/datasets/dayveedd/gsm8k-lowres-math}{\underline{gsm8k-lowres-math dataset}}.
    \item Benchmarked Qwen2.5-1.5B-Instruct and Tiny Aya families of models on math tasks across all three languages to guide model selection before fine-tuning.
\end{itemize}
\vspace{6pt}
\project{ECG Benchmarking for Atrial Fibrillation Detection}{PySpark, PostgreSQL, MinIO}
\begin{itemize}
    \item Contributed to a reproducible benchmarking framework spanning 42+ ECG databases, with deterministic data splits, cached feature generation, and traceable experiment artifacts.
    \item Helped benchmark 14+ detectors across datasets using window- and recording-level metrics, with quality and runtime comparisons through configurable experiments.
\end{itemize}
\newpage
\section{Selected Projects (continued)}
\project{Kinship Recognition from Photographs}{ArcFace, SentenceTransformers, FAISS, CatBoost}
\begin{itemize}
    \item Contributed to dual-modal face/name retrieval on Families in the Wild, combining ArcFace and MiniLM embeddings through two FAISS indexes and reciprocal rank fusion.
    \item Worked on CatBoost classifiers for family membership and relationship type across 11 kinship categories, with SHAP-based explanations.
\end{itemize}
\vspace{6pt}
\project{RAG for Document Collections}{Python, SentenceTransformers, LLMs}
\begin{itemize}
    \item Contributed to building a Wikipedia-based RAG pipeline with article cleaning, chunking, 384-dimensional MiniLM embeddings, and incremental embedding updates.
    \item Combined 3-gram candidate retrieval with cosine-similarity reranking to supply relevant document snippets for LLM-generated answers.
\end{itemize}

\section{Research}
\textbf{Meta-learners and LLM-Guided Adaptive Weighting for Physics-Informed Neural Networks}\par
\begin{itemize}
    \item Lead-authored a study introducing an LLM meta-controller that dynamically balances PDE-residual and boundary-condition losses using training history, validation trends, and stability diagnostics.
    \item Benchmarked LLM-guided weighting against fixed weights and LBPINN on Poiseuille flow across 100 pressure-gradient values and five random seeds, assessing convergence, solution accuracy, and adaptation trade-offs.
\end{itemize}
\vspace{6pt}
\textbf{Recommendation Systems via Approximate Matrix Factorization}\par
\begin{itemize}
    \item Co-authored a project report comparing six matrix-factorization methods for collaborative filtering, evaluating RMSE, convergence, and sensitivity to rank and regularization.
\end{itemize}
\vspace{6pt}
\textbf{Software Engineering Metrics Selection Using Equilibrium Optimizer and Differential Evolution}\par
\begin{itemize}
    \item Implemented the Equilibrium Optimizer for software-metric selection, using heuristic search to identify compact feature subsets while preserving geometric relationships between data points.
    \item Used Sammon error and Kruskal stress to evaluate structural preservation; the Equilibrium Optimizer reduced 30 metrics to five in a co-authored study also evaluating Differential Evolution.
\end{itemize}

\end{document}